\documentclass[conference]{IEEEtran}
\IEEEoverridecommandlockouts

\usepackage[T1]{fontenc}
\usepackage[utf8]{inputenc}
\usepackage{microtype}
\usepackage{amsmath,amssymb}
\usepackage{booktabs}
\usepackage{tabularx}
\usepackage{array}
\usepackage{enumitem}
\usepackage{xcolor}
\usepackage{graphicx}
\usepackage{hyperref}

\definecolor{accent}{HTML}{B11226}
\definecolor{todocol}{HTML}{B11226}
\hypersetup{colorlinks=true, linkcolor=accent, citecolor=accent, urlcolor=accent,
  pdftitle={Diffusion-Trained Backbones for Flow-Matching VLA Policies},
  pdfauthor={Ryan Rahman}}

% Living-document marker: anything still to fill shows in red.
\newcommand{\todo}[1]{\textcolor{todocol}{\textbf{[TODO: #1]}}}
\newcommand{\tbd}{\textcolor{todocol}{TBD}}

\title{Diffusion-Trained Multimodal Backbones for Action-Token
Representations in Flow-Matching Vision-Language-Action Policies}
\author{\IEEEauthorblockN{Ryan Rahman}
\IEEEauthorblockA{University of Waterloo}
\thanks{Working preprint; updated as experiments complete.}}

\begin{document}
\maketitle

\begin{abstract}
Vision-Language-Action (VLA) policies pair a pretrained multimodal backbone with
a continuous action generator such as a flow-matching action expert. Most
backbones inherit an autoregressive (AR) next-token pretraining prior, yet robot
action chunks are block-structured rather than left-to-right. We ask whether a
\emph{diffusion-trained} bidirectional backbone produces better action-token
hidden states than a scale-matched AR backbone when both drive the same
flow-matching action expert. We isolate the pretraining prior with a matched pair
--- Dream-7B (discrete diffusion, initialized from Qwen2.5-7B) versus Qwen2.5-7B
(AR) --- holding architecture, initialization, action expert, dataset, and control
stack fixed, and evaluate entirely in simulation on the LIBERO benchmark with
both policy-level success and representation-level probing. On the matched $C-D$
contrast (single seed), the diffusion-trained backbone improves mean LIBERO
success by $+5.9$ points (82.2 vs.\ 76.3), with the largest gains on the
long-horizon (LIBERO-10, $+12.0$) and object (\,$+11.2$) suites --- precisely the
regimes where the autoregressive backbone is weakest --- and parity on the
near-ceiling spatial suite. Yet linear probing of the two backbones' action-token
features finds them equally decodable for low-level action, state, and gripper
targets ($|C-D|\le0.01$): a rollout--probe dissociation indicating the diffusion
prior's benefit is not reducible to the linear decodability of immediate control
signals, but acts through higher-level or closed-loop structure.
\end{abstract}

\begin{IEEEkeywords}
vision-language-action policies, flow matching, diffusion models, imitation
learning, representation probing
\end{IEEEkeywords}

\section{Introduction}
State-of-the-art vision-language-action (VLA) policies pair a pretrained
vision-language model (VLM) backbone with a flow-matching action expert that turns
the backbone's hidden states into continuous action chunks, as in $\pi_0$ and
$\pi_{0.5}$ \cite{pi0,pi05,openpi}. The action expert generates a whole block of
future actions at once, refining them iteratively and attending across the block
bidirectionally --- a generative process much closer to denoising than to
left-to-right prediction. The backbone that feeds it, however, is almost always an
\emph{autoregressive} VLM, pretrained to predict the next token causally. This
raises a question that the field has largely left implicit: does the backbone's
\emph{pretraining objective} --- the prior it brings to the action-token hidden
states --- matter for the policy, and would a backbone pretrained by denoising
suit a flow-matching expert better than a matched autoregressive one?

The question has been hard to ask cleanly for two reasons. First, until recently
there was no diffusion-pretrained language model at the scale of the autoregressive
backbones used in VLAs; discrete-diffusion LMs such as Dream-7B
\cite{dream}, initialized from Qwen2.5-7B \cite{qwen25}, now provide an
architecturally identical counterpart that differs \emph{only} in pretraining
objective. Second, the inference-time attention mask is easy to confuse with the
pretraining prior: we note that the $\pi_{0.5}$ backbone is in fact a prefix-LM
whose action block is already attended bidirectionally, so the causal-vs-bidirectional
distinction at inference is not the variable of interest --- the pretraining prior
is. These two observations make a controlled comparison possible: hold the action
expert, vision encoder, data, and adaptation recipe fixed, and swap only the
backbone between a diffusion-trained model (Dream-7B) and its autoregressive
initializer (Qwen2.5-7B).

We run exactly this comparison on LIBERO. The diffusion-trained backbone yields
the better policy --- $+5.9$ points of mean success and $+12.0$ on the
long-horizon suite, precisely where the autoregressive backbone is weakest --- at
matched scale, initialization, and everything downstream of the backbone, so the
gap is attributable to the pretraining prior. Probing the two backbones'
action-token features, however, finds them equally linearly decodable for
low-level control targets: the behavioral advantage is real but does not reduce to
the linear readability of immediate actions, a dissociation that locates the
effect in higher-level or closed-loop structure.

Our contributions are: (i) a scale/architecture-matched isolation of the
diffusion pretraining prior for VLA backbones; (ii) a precise characterization of
the \(\pi_{0.5}\) backbone as a prefix-LM (bidirectional action block), separating
attention-mask topology from the pretraining prior; (iii) representation-level
probing of action-token hidden states; (iv) a controlled LIBERO evaluation
directly comparable to published diffusion-VLA results.

\section{Related Work}

\paragraph{Diffusion and flow policies.} Generating actions by denoising or flow
matching is now standard in imitation learning: Diffusion Policy \cite{diffusion_policy}
and 3D Diffusion Policy \cite{dp3} show that iterative, block-structured action
generation captures multimodal behavior better than regressing a single action.
These methods build the denoiser on task-specific or from-scratch encoders; we
inherit the flow-matching action head but ask what its \emph{backbone} should be.

\paragraph{VLM backbones with flow-matching experts.} $\pi_0$ and $\pi_{0.5}$
\cite{pi0,pi05,openpi} scale this recipe by conditioning a flow-matching expert on
a large pretrained VLM, achieving strong generalist manipulation. In these systems
the backbone is autoregressively pretrained and (in $\pi_{0.5}$) used as a
prefix-LM, so the action block is already bidirectional at inference; the
pretraining objective of the backbone is not itself treated as a design choice.
Our study takes that objective as the variable.

\paragraph{Diffusion-trained and multimodal-diffusion LMs.} Discrete-diffusion
language models --- Dream-7B \cite{dream}, and encoder-decoder diffusion LMs such
as DiffusionGemma \cite{diffusiongemma_docs} --- and multimodal diffusion models
like LaViDa \cite{lavida} establish that denoising-pretrained backbones exist at
VLM scale. Crucially, Dream is initialized from an autoregressive model
(Qwen2.5-7B), giving us a matched pair whose only systematic difference is the
pretraining objective.

\paragraph{Unified-backbone diffusion VLAs (closest work).} A recent line folds
action decoding \emph{into} a diffusion backbone --- Dream-VLA \cite{dreamvla},
LLaDA-VLA \cite{lladavla}, and Discrete Diffusion VLA \cite{discretediffvla} ---
reporting strong LIBERO results. These works differ from ours in three ways that
matter for what can be concluded: they (a) do not hold a matched autoregressive
backbone fixed as a control, so the effect of the pretraining prior is
confounded with architecture and training changes; (b) let the backbone generate
actions directly rather than through a separate, fixed flow-matching expert, so
representation quality and decoding are entangled; and (c) do not probe the
learned representations. We instead keep one external flow-matching expert
constant, swap only the backbone against its own AR initializer, and add
representation-level probing --- trading peak performance for a controlled,
interpretable contrast.

\paragraph{Representation quality and action-token structure.} Crossway Diffusion
\cite{crossway} shows that auxiliary representation learning improves diffusion
policies, and work on ordered/structured action tokens \cite{oat} highlights that
how actions are laid out as tokens affects learning --- both motivating our focus
on the \emph{action-token hidden states} the expert consumes, and our attempt to
measure their quality directly via probing rather than only through rollouts.

\section{Method}
\subsection{Backbone-swap with a shared flow-matching expert}
Image/language/proprio/flow-time/noisy-action tokens are processed by the
backbone, which contextualizes the action tokens; the hidden states at the
action-token positions, $H_{\text{action}}$, are read by a flow-matching action
expert that outputs the vector field over the action chunk. The flow objective
is $\mathcal{L}=\lVert v_\theta - (A_{\text{noise}}-A_{\text{clean}})\rVert^2$
with $x_t = t\,A_{\text{noise}} + (1-t)\,A_{\text{clean}}$.

\paragraph{Standalone vs.\ fused expert.} A subtlety determines the design.
$\pi_{0.5}$'s action expert is not a detachable head: it is a second transformer
stream \emph{fused} into the backbone's self-attention (per layer, backbone and
expert queries/keys/values are concatenated and attended jointly, sharing RoPE
and KV cache), which forces the expert to match the backbone's layer count, head
geometry, and normalization. This fusion cannot be preserved when the backbone
changes from Gemma to Qwen2.5/Dream (different depth, heads, RoPE, biases). We
therefore use a \emph{standalone} flow expert --- a small DiT-style transformer
over the action tokens, flow-time injected via adaLN-Zero conditioning --- that
reads $H_{\text{action}}$ and is decoupled from the backbone's attention. It is
shared \emph{identically} across variants A$'$/B/C/D; the only variant-specific
parameter is a linear input adapter mapping the backbone width to the common
expert width ($\sim$2--4M of $\sim$117M parameters), so the C$-$D contrast
remains clean. Model A (stock $\pi_{0.5}$, fused expert) is retained as a
reference that reproduces published numbers, and A$'$ (Gemma backbone $+$ the
standalone expert) is the architecture-matched baseline.

\subsection{Model variants}
The action expert, action space, flow objective, dataset, and control stack are
identical across all rows; only the backbone/attention topology changes.

\begin{table*}[t]
\centering
\caption{Model variants. Identifying contrast: $C-D$ (diffusion prior at matched
scale/architecture).}
\begin{tabularx}{\linewidth}{>{\bfseries}p{0.06\linewidth} X}
\toprule
 & Backbone \\ \midrule
A & Stock $\pi_{0.5}$ (Gemma, \emph{fused} expert) --- reference baseline \\
A$'$ & Gemma backbone $+$ standalone expert --- architecture-matched baseline \\
D & Qwen2.5-7B (AR) $+$ standalone expert --- matched control \\
C & Dream-7B (diffusion, init.\ from Qwen2.5-7B) $+$ standalone expert --- proposed \\
B & Model D's backbone $+$ expanded attention topology --- topology ablation \\
\bottomrule
\end{tabularx}
\end{table*}

\subsection{Reading hidden states from a diffusion backbone}
\label{sec:kread}
A diffusion-trained backbone could in principle be run for several internal
refinement steps ($k>1$) before its hidden states are read, unrolling the
denoising process the AR backbone lacks. To keep the $C-D$ contrast about the
pretraining \emph{prior} and not about extra inference compute, we read both
backbones with a single forward pass ($k{=}1$): C and D then do exactly the same
amount of work per action, and any gap is due to the learned representation, not
to additional refinement. This backbone read is also distinct from --- and
happens once per --- each flow-matching step of the standalone expert. We report
$k{=}1$ throughout; multi-step refinement reads ($k\in\{2,4\}$), which trade
compute for a potentially richer read, are a natural ablation left to future work.

\section{Experimental Setup}
Evaluation is entirely in simulation on LIBERO \cite{libero} (single-arm Franka,
7-DoF), across four suites (Spatial, Object, Goal, Long/libero\_10), 10 tasks per
suite $\times$ 50 trials $=$ 500 rollouts/suite. Backbones are $\sim$7B (Dream-7B
\cite{dream} / Qwen2.5-7B \cite{qwen25}); Model A is the stock $\pi_{0.5}$ (Gemma)
backbone. \textbf{Vision is held identical across the standalone-expert variants}
(A$'$/C/D): the SigLIP tower from PaliGemma plus a fresh linear connector to the
backbone width, so only the language backbone changes between C and D. Training
forks the openpi \texttt{pi05\_libero} recipe (30k steps, batch 32, LR $5\times
10^{-5}$, action horizon $H{=}10$, action dim padded to 32). \textbf{Adaptation
differs by variant and is stated with each result:} A and A$'$ are fully
fine-tuned from \texttt{pi05\_base}; C and D use an identical strengthened-LoRA
recipe (rank 64, $\alpha{=}128$ on the frozen 7B backbone, with the SigLIP tower,
connector, input adapters, and standalone expert trainable) --- chosen because
full fine-tuning of a 7B backbone with AdamW does not fit one 80GB GPU under DDP,
and because freezing the base best preserves the pretraining prior the $C-D$
contrast is designed to isolate. The diffusion backbone is read with a single
forward pass ($k{=}1$), matching D exactly; deeper $k$-step refinement reads are
left to future work. Each training run used 2$\times$A100 ($\sim$5\,s/step). We
report per-suite success rate, with LIBERO-Long as the primary discriminator,
plus representation probes (\S\ref{sec:probing}). A and A$'$ are averaged over 3
seeds; C and D are single-seed ($n{=}1$) --- the identifying $C-D$ contrast is
evaluated on the matched seed, with additional seeds for error bars left to
future work.

\section{Results}

\subsection{Stage 1 --- Model A baseline reproduction}
Our Model A reproduces the published $\pi_{0.5}$ LIBERO results to within
$\sim$1\% on every suite (average 96.6 vs.\ 96.85), validating the training and
evaluation harness before any backbone swap. This holds despite training at
batch size 32 (vs.\ the published 256), supporting our batch choice. On the
long-horizon suite Model A slightly exceeds the published number (93.1 vs.\ 92.4).

\begin{table*}[t]
\centering
\caption{Model A (stock $\pi_{0.5}$) LIBERO success rate (\%), mean $\pm$ sd over
3 seeds, vs published $\pi_{0.5}$@30k. Higher is better.}
\begin{tabular}{lccccc}
\toprule
Model & Spatial & Object & Goal & Long & Avg \\
\midrule
Published $\pi_{0.5}$@30k & 98.8 & 98.2 & 98.0 & 92.4 & 96.85 \\
A (ours, 3 seeds) & 97.7\,$\pm$\,0.6 & 98.3\,$\pm$\,0.6 & 97.1\,$\pm$\,0.5 & 93.1\,$\pm$\,1.4 & 96.6\,$\pm$\,0.2 \\
\bottomrule
\end{tabular}
\end{table*}

\subsection{Main comparison: diffusion vs matched AR backbone}
\label{sec:main}
\begin{table*}[t]
\centering
\caption{LIBERO success rate (\%) by variant. A/A$'$ are mean over 3 seeds;
D/C are single-seed (n=1, seed 0) preliminary runs (error bars pending). The
identifying contrast is $C-D$, evaluated on the matched single seed. D/C use the
strengthened LoRA recipe (unfrozen SigLIP, rank 64); A$'$ is full fine-tuned.}
\begin{tabular}{lccccc}
\toprule
Model & Spatial & Object & Goal & Long & Avg \\
\midrule
A (stock $\pi_{0.5}$, fused, 3 seeds) & 97.7 & 98.3 & 97.1 & 93.1 & 96.6 \\
A$'$ (Gemma + standalone, 3 seeds)    & 96.7 & 92.6 & 90.5 & 79.6 & 89.9 \\
D (Qwen2.5-7B, AR, n{=}1)             & 94.2 & 82.0 & 71.4 & 57.4 & 76.3 \\
C (Dream-7B, diff., n{=}1)   & 93.0 & 93.2 & 73.2 & 69.4 & 82.2 \\
B (topology)                 & \tbd & \tbd & \tbd & \tbd & \tbd \\
\midrule
$C-D$ (diffusion prior)      & $-1.2$ & $+11.2$ & $+1.8$ & $+12.0$ & $+5.9$ \\
$A'\!-\!A$ (decoupling cost) & $-1.0$ & $-5.7$ & $-6.6$ & $-13.5$ & $-6.7$ \\
\bottomrule
\end{tabular}
\end{table*}

\paragraph{The diffusion prior helps, and most where AR is weakest.} On the
identifying $C-D$ contrast (matched seed, identical standalone expert, LoRA
recipe, and vision), the diffusion-trained backbone improves mean success by
$+5.9$ points. The gain is not uniform: it concentrates on \textbf{long-horizon}
(LIBERO-10, $+12.0$) and \textbf{object} ($+11.2$), is small on goal ($+1.8$), and
is a statistical tie on the near-ceiling spatial suite ($-1.2$). This is the
predicted pattern --- D (and A$'$) degrade sharply on long, multi-stage tasks
(A$'-$A is $-13.5$ on long), and it is exactly there that the bidirectional,
diffusion-trained representation recovers the most. The effect is measured at
matched scale (7B), matched initialization (Dream is initialized from Qwen2.5-7B),
and matched everything-else, so it is attributable to the pretraining objective
rather than capacity or architecture. \emph{Caveat:} $C$ and $D$ are single-seed
($n{=}1$); seeds 1--2 are pending for error bars, and the spatial cell for $C$ is
recovered from clean per-task rates after an eval-harness logging fault (see
RESULTS.md).

\subsection{Representation-level probing}
\label{sec:probing}

The rollout comparison (\S\ref{sec:main}) tells us \emph{whether} the diffusion
prior helps end-to-end; probing asks \emph{why}, by measuring how much
action-relevant information is already linearly decodable from the backbone's
action-token hidden states before the flow expert acts. Because C and D share
architecture, width, vision, LoRA recipe, and expert---differing only in
pretraining objective---a probing gap between them is attributable to the
diffusion prior, and mirrors the $C-D$ contrast at the representation level.

\paragraph{Feature extraction.}
For each model we run a single forward pass over a held-out set of LIBERO
observations and read the last-layer backbone hidden states at the action-token
positions, $h \in \mathbb{R}^{H\times d}$ ($H{=}32$ action steps, $d{=}3584$ for
Qwen/Dream, $2048$ for Gemma). The backbone is \emph{frozen}: no gradients flow
into it, and the flow expert is not involved. To get a fixed-length feature we
report both mean-pooling over the $H$ action tokens and the per-token features
(probing each step $t$ independently), so decodability can be traced along the
chunk. Features are extracted at flow time $\tau{=}1$ (pure-noise action input),
so the probe sees only what the \emph{conditioning} carries, not a partially
denoised action; for the diffusion backbone we use the $k{=}1$ single-pass read
(\S\ref{sec:kread}) to match D exactly.

\paragraph{Probe targets.} We decode a battery of action-relevant quantities of
increasing abstraction, each with its own natural metric:
\begin{itemize}
  \item \textbf{Next action / chunk} (regression, $R^2$): the ground-truth
        end-effector delta at step $t$, and the full 32-step chunk. Direct test
        of how much of the control signal is linear in the features.
  \item \textbf{Gripper state} (binary, accuracy/AUC): open vs.\ closed, and the
        \emph{transition} step (the contact/grasp event), which is where
        long-horizon tasks fail.
  \item \textbf{Subtask phase} (multi-class, accuracy): the current phase within
        a multi-stage task (reach / grasp / transport / place), the abstraction
        most relevant to the Long suite where $C-D$ is hypothesized to be largest.
  \item \textbf{Target-object pose} (regression, $R^2$): position of the task's
        target object, testing spatial grounding of the visual conditioning.
  \item \textbf{Chunk success} (binary, AUC): whether the executed chunk advances
        the task, giving a representation-level analogue of rollout success.
\end{itemize}

\paragraph{Probe model and protocol.} To measure \emph{linear} decodability we
use a linear (ridge / logistic) probe; we additionally report a small 2-layer MLP
probe as a nonlinear ceiling. Probes are \textbf{capacity-matched} across
backbones (identical probe architecture and hyperparameters; the only input
difference is $d$, and we control for it by also reporting probes on a
PCA-reduced common dimensionality). We fit on a train split and report on a held-out
test split, with 5 seeds for the probe fit (backbone frozen) to get error bars on
the probe metrics themselves. Chance and a shuffled-label baseline are reported
per target.

\paragraph{Two settings.}
\begin{enumerate}
  \item \textbf{Trained checkpoints (primary).} Probe the LoRA-adapted C and D
        backbones---the exact weights that produced the rollouts in
        \S\ref{sec:main}. This ties the probe gap to the success gap on the same
        models. This is the reported comparison.
  \item \textbf{Stock backbones (isolation).} Probe the off-the-shelf Gemma,
        Qwen2.5-7B, and Dream-7B backbones with \emph{no} VLA fine-tuning
        (vision features held fixed across the three). This needs no policy
        checkpoint---including no A$'$ checkpoint---and isolates the pretraining
        prior in its purest form: three pretraining objectives, everything else
        held constant. A diffusion advantage that is already present here is
        evidence the prior, not the fine-tuning, drives it.
\end{enumerate}

\paragraph{Reading the result.} The key comparison is probe metric on C vs.\ D
per target. Three outcomes and their reading: (i) \emph{C probes higher and rolls
out higher} --- the diffusion prior yields more decodable action representations,
and that translates to control (the clean positive result); (ii) \emph{C probes
higher but does not roll out higher} --- a probe-win/rollout-win
\textbf{dissociation}, meaning the information is present but the fixed expert
cannot exploit it (implicating the decoupled-expert design, cf.\ $A'\!-\!A$);
(iii) \emph{no probe gap} --- the objectives produce equivalent action-token
representations at this scale. We expect the effect, if present, to be largest on
the phase and gripper-transition targets, matching the Long-suite regime.
\paragraph{Implementation.} The pipeline is
\texttt{scripts/extract\_probe\_features.py} (dumps the frozen $\tau{=}1$
action-token features from a trained C/D checkpoint, mean-pooled and first-token,
alongside labels) and \texttt{openpi.diffusion\_backbone.probing} (fits the
capacity-matched linear + MLP probes over 5 splits and emits the C-vs-D table).
C and D are extracted in the same dataset order so they are probed on identical
frames; the probe step asserts this alignment. The first implemented target set is
the one derivable directly from the dataset --- immediate action and full action
chunk (real 7-DoF, $R^2$), proprioceptive state ($R^2$), and gripper command
(median-split, accuracy/AUROC). The more abstract targets (subtask phase,
target-object pose, chunk success) require simulator-state annotation and are
deferred to a labeling pass.

\begin{table}[t]
\centering
\caption{Linear-probe predictivity of frozen action-token features (mean-pooled,
$\tau{=}1$), C vs.\ D, seed 0, 4000 matched frames, 5 splits. $R^2$ for
regression, accuracy for gripper. Both backbones are probed on \emph{identical}
frames. The MLP (nonlinear) probe did not converge (max-iter reached, high
variance) and is omitted as inconclusive.}
\label{tab:probing}
\begin{tabular}{lccc}
\toprule
Target & C (Dream) & D (Qwen) & $C-D$ \\
\midrule
action (next step), $R^2$      & 0.874 & 0.869 & $+0.006$ \\
action (full chunk), $R^2$     & 0.889 & 0.887 & $+0.002$ \\
proprioceptive state, $R^2$    & 0.964 & 0.954 & $+0.009$ \\
gripper (median-split), acc    & 0.969 & 0.969 & $+0.001$ \\
\bottomrule
\end{tabular}
\end{table}

\paragraph{A rollout--probe dissociation.} At the representation level, C and D
are \emph{tied}: their action-token features are equally --- and very highly ---
linearly decodable for every low-level target ($|C-D|\le0.01$ throughout). This
does not track the rollout result, where $C-D$ is $+5.9$ on average and $+12.0$ on
long-horizon (\S\ref{sec:main}). The behavioral advantage of the diffusion prior
is therefore \emph{not} explained by better linear decodability of low-level
action, state, or gripper information --- that information is equally present in
both backbones.

Three factors frame this null. (i) \textbf{Ceiling:} both backbones probe at
$0.87$--$0.96$, leaving little headroom for a linear probe to separate them, so the
result is ``no difference detectable on these targets,'' not ``identical
representations.'' (ii) \textbf{Target level:} these targets are low-level and
read from single, in-distribution frames, whereas the rollout gap concentrates on
long-horizon control --- a closed-loop, compounding-error regime that static
single-frame probing cannot capture; the abstract targets most likely to separate
the backbones (subtask phase, spatial relations) are the ones deferred for
requiring simulator-state labels. (iii) \textbf{Nonlinear probe:} the MLP did not
converge and is inconclusive. The takeaway is that the diffusion prior's benefit
appears to act through higher-level/temporal structure or closed-loop robustness
rather than through the linear decodability of immediate control signals --- ruling
out the simplest mechanistic explanation and motivating the deferred
abstract-target and closed-loop probes.

\section{Discussion}

The matched $C-D$ contrast lands on the \emph{positive} outcome of the three we
pre-registered: the diffusion-trained backbone yields better action-token
representations, and the advantage survives end-to-end into closed-loop control
($+5.9$ mean, $+12.0$ on long-horizon). Because $C$ and $D$ are identical up to
the pretraining objective (Dream is initialized from Qwen2.5-7B; same expert,
vision, LoRA recipe, seed, and data), we read this as evidence that a
bidirectional/denoising pretraining prior produces hidden states that a fixed
flow-matching expert can exploit better than those of a matched autoregressive
backbone --- most in the long-horizon, multi-stage regime where the AR backbone
degrades.

The probing analysis (\S\ref{sec:probing}) sharpens \emph{how} this happens by
ruling out the simplest account. If the diffusion prior simply made low-level
control more linearly readable, C would out-probe D; instead the two are tied at
ceiling on every low-level target (Table~\ref{tab:probing}) while C clearly
out-\emph{rolls} D. This rollout--probe dissociation says the advantage does not
live in the linear decodability of immediate action/state signals --- both
backbones carry those equally --- but in something the static, single-frame,
in-distribution probe does not see: plausibly higher-level or temporal structure,
or closed-loop robustness to compounding error, consistent with the gap being
largest on long-horizon tasks. Discriminating these requires the deferred
abstract-target probes (subtask phase, spatial relations) and a closed-loop /
distribution-shift probe, which we leave to future work.

\todo{once B is in: B$-$D (topology) vs C$-$D (prior). Note C sits below A$'$/A in
absolute terms (decoupling + LoRA-vs-full-FT cost, common to C and D, cancels in
C$-$D).}

\paragraph{Limitations.} $C$/$D$ are single-seed ($n{=}1$); the $+5.9$/$+12.0$
gaps are large relative to the A-seed spread ($\pm0.2$ avg) but seeds 1--2 are
needed for formal error bars. Evaluation is simulation-only (LIBERO), so external
validity to real robots is untested. The absolute numbers reflect the standalone
(decoupled) expert, which costs $\sim$6--7 points vs.\ the fused expert
($A'-A$); this cost is shared by $C$ and $D$ and thus does not affect the
identifying contrast.

\section{Conclusion}

We asked whether the pretraining objective of a VLA backbone --- diffusion versus
autoregressive --- matters for downstream manipulation when everything else is
held fixed, and built a controlled test to answer it: a single standalone
flow-matching action expert driven by interchangeable 7B backbones that differ
only in how they were pretrained (Dream, initialized from Qwen2.5-7B, versus
Qwen2.5-7B itself). The diffusion-trained backbone is the better policy, improving
mean LIBERO success by $+5.9$ points and by $+12.0$ on the long-horizon suite ---
the regime where the autoregressive backbone, and the decoupled expert generally,
degrade most. Because the two systems are matched in scale, initialization,
vision, action expert, adaptation recipe, data, and seed, we attribute the gap to
the pretraining prior rather than to capacity or architecture. At the same time,
linear probing shows the two backbones' action-token representations to be equally
decodable for low-level action, state, and gripper targets: the behavioral
advantage does not reduce to the linear readability of immediate control signals,
a dissociation that points instead toward higher-level or closed-loop structure.
The study is single-seed, simulation-only, and probes a deliberately low-level
target set, so we read it as a controlled existence proof rather than a final
measurement: a diffusion-trained backbone can yield a better flow-matching policy
than a matched autoregressive one, most where it matters, and not for the reason
one would first guess. Establishing the mechanism --- through abstract-target and
closed-loop probing --- and the generality --- through additional seeds, the
attention-topology control ($B$), and real-robot evaluation --- is the natural
next step.

\bibliographystyle{IEEEtran}
\begin{thebibliography}{99}
\bibitem{diffusion_policy} Chi et al. \textit{Diffusion Policy}. RSS 2023. \url{https://arxiv.org/abs/2303.04137}
\bibitem{pi0} Physical Intelligence. \textit{$\pi_0$}. 2024. \url{https://www.physicalintelligence.company/download/pi0.pdf}
\bibitem{pi05} Black et al. \textit{$\pi_{0.5}$}. CoRL 2025. \url{https://proceedings.mlr.press/v305/black25a.html}
\bibitem{openpi} Physical Intelligence. \textit{OpenPI}. \url{https://github.com/Physical-Intelligence/openpi}
\bibitem{dream} Ye, Xie et al. \textit{Dream 7B}. 2025. \url{https://arxiv.org/abs/2508.15487}
\bibitem{qwen25} Qwen Team. \textit{Qwen2.5 Technical Report}. 2024. \url{https://arxiv.org/abs/2412.15115}
\bibitem{lavida} Li et al. \textit{LaViDa}. 2025. \url{https://arxiv.org/abs/2505.16839}
\bibitem{dreamvla} DreamLM. \textit{Dream-VL / Dream-VLA}. 2025--2026. \url{https://hkunlp.github.io/blog/2025/dream-vlx/}
\bibitem{lladavla} Wen et al. \textit{LLaDA-VLA}. 2025. \url{https://wenyuqing.github.io/llada-vla/}
\bibitem{discretediffvla} \textit{Discrete Diffusion VLA}. 2025. \url{https://arxiv.org/abs/2508.20072}
\bibitem{diffusiongemma_docs} Google. \textit{DiffusionGemma}. 2026. \url{https://ai.google.dev/gemma/docs/diffusiongemma}
\bibitem{crossway} Li et al. \textit{Crossway Diffusion}. ICRA 2024. \url{https://arxiv.org/abs/2307.01849}
\bibitem{oat} Liu et al. \textit{Ordered Action Tokens}. 2026. \url{https://arxiv.org/abs/2607.21670}
\bibitem{dp3} Ze et al. \textit{3D Diffusion Policy}. 2024. \url{https://arxiv.org/abs/2403.03954}
\bibitem{libero} Liu et al. \textit{LIBERO}. NeurIPS D\&B 2023. \url{https://arxiv.org/abs/2306.03310}
\end{thebibliography}

\end{document}
